% Rectificación quirúrgica: generación conjunta, profundidad arbitraria y
% estatuto tipado de los transportes ternarios.  Este bloque pertenece al
% capítulo 27 y precede a la separación expositiva de los cuatro lectores.

\section{Théorème conjoint de génération coinductive et de profondeur arbitraire}
\label{sec:c27-generacion-coinductiva-profundidad-arbitraria}

La structure de départ de ce théorème n'est pas l'ensemble nu de neuf
symboles. C'est le système typé
\[
 \mathfrak H_9=
 \bigl(\mathrm{APP},\mathrm{TRIT},\mathrm{TPK};
       \widehat{\mathcal X}^{\mathrm{enr}},\mathcal U,
       \Gamma_9,\operatorname{Ext},\operatorname{tr}\bigr),
\]
construit sur la double réalisation positionnelle et arithmétique de
\(\mathcal D_9=\{1,\ldots,9\}\). APP accumule et replie au moyen de ses feuillets
\(\Sigma\) et \(\Pi\), en conservant résidu et quotient ; TRIT oriente chaque
transition dans l'un de ses trois régimes ; et TPK sélectionne, transporte,
mémorise et revient sur la fibre enrichie. En particulier,
\[
 M_{\rm ph}:\mathcal D_9\longrightarrow\{+1,-1,0\},
 \qquad
 m_{\rm ph}:=\bigl(M_{\rm ph}(1),\ldots,M_{\rm ph}(9)\bigr)^{\mathsf T},
\]
\[
 \operatorname{Tra}_t(y,d):=\operatorname{Lift}(T_d)(y),
 \qquad
 \operatorname{Upd}_t:=\operatorname{Rec}_t\circ\operatorname{Mem}_t,
\]
et, pour tout \(x\in\widehat{\mathcal X}^{\mathrm{enr}}\),
\[
 \mathcal U_t(x)
 =\operatorname{Upd}_t\!\left(
   \operatorname{Tra}_t\!\left(
    \operatorname{Sel}_t\!\left(x,M_{\rm ph}(g(t))\right),d(t)
   \right)\right).
\]
Cette composition est la donnée dynamique effective : la carte à neuf symboles
est son support fini, non un substitut de la dynamique. En particulier,
\(M_{\rm ph}(g(t))\) est la valeur scalaire que reçoit le sélecteur, tandis
que \(m_{\rm ph}\) est le vecteur des neuf valeurs de la fonction ; aucun
ne se compose directement avec un opérateur d'état.

Pour \(r\in6\mathbb N_{>0}\), soient
\(\widehat{\mathcal X}^{\mathrm{enr}}_r\) les états tronqués avec feuillet,
orientation, résidu, quotient, retenue, voie, bord et mémoire, et soient
\[
 \operatorname{Ext}_{r+6,r}\subseteq
 \widehat{\mathcal X}^{\mathrm{enr}}_r\times
 \widehat{\mathcal X}^{\mathrm{enr}}_{r+6},
 \qquad
 \operatorname{tr}_{r+6,r}:
 \widehat{\mathcal X}^{\mathrm{enr}}_{r+6}
 \longrightarrow
 \widehat{\mathcal X}^{\mathrm{enr}}_r
\]
la relation de prolongement admissible et la troncature. Pour
\(x_r\in\widehat{\mathcal X}^{\mathrm{enr}}_r\), la fibre locale
\[
 \operatorname{Ext}_{r+6,r}(x_r)
 :=\{x_{r+6}:(x_r,x_{r+6})\in\operatorname{Ext}_{r+6,r}\}
\]
peut contenir plusieurs enfants. La compatibilité correcte est
\[
 (x_r,x_{r+6})\in\operatorname{Ext}_{r+6,r}
 \Longrightarrow
 \operatorname{tr}_{r+6,r}(x_{r+6})=x_r.
\]
Dans \(R_{36}\), les trois lignes retrouvées
\(222220,021222,102011\), transportées au calibre actif par la permutation
\(p_0=(0,1,2,4,5,3)\) en indices commençant à zéro ---de manière équivalente,
\(p_1=(1,2,3,5,6,4)\) en indices commençant à un---, portent avant toute
réalisation réelle les caractères
internes de clôture, de propagation et d'auto-échelle. Leur triplet, avec le
registre signé \(U_{12}\) et le sceau \(K\), porte la coordonnée interne de
\(\alpha\). Pour chaque
\(\chi\in\{\pi,\varphi,e,\alpha\}\), nous écrivons l'état ainsi enrichi
comme
\[
 x_r^\chi=(\widetilde x_r,\chi_{\rm HMT},
            \varepsilon_r^\chi,N_r^\chi).
\]
Ici \(\varepsilon_r^\chi\) est le résidu interne produit par le quotient
APP dans l'état, non \(\operatorname{frac}(P_{\rm ar}(\chi))\). La division
euclidienne interne fixe
\[
 d_r^\chi=\operatorname{Dig}_{729}(\varepsilon_r^\chi),
 \quad 0\le d_r^\chi<729,
 \quad
 \varepsilon_{r+6}^\chi=729\varepsilon_r^\chi-d_r^\chi,
 \quad
 N_{r+6}^\chi=729N_r^\chi+d_r^\chi.
\]
La sélection, le transport et la mise à jour sur la section
\(\chi\) sont typés par
\[
 \operatorname{Sel}_r^\chi(x)
 :=\operatorname{Sel}\!\left(
 x,M_{\rm ph}(\rho_9(r/6+1))\right),
 \qquad
 \operatorname{Tra}_r^\chi(y,d):=\operatorname{Lift}(T_d)(y),
 \qquad
 \operatorname{Upd}_r^\chi:=\operatorname{Rec}_r^\chi\circ
 \operatorname{Mem}_r^\chi.
\]
La transition complète est
\[
 \mathcal U_r^\chi(x)
 :=\operatorname{Upd}_r^\chi\!\left(
  \operatorname{Tra}_r^\chi\!\left(
   \operatorname{Sel}_r^\chi(x),d_r^\chi
  \right)\right),
\]
et ne lit que l'état enrichi précédent : phase, mémoire, retenue, feuillet,
orientation, voie, bord, signature, registre, caractère et résidu internes.
Elle ne lit pas un chiffre futur, \(P_{\rm ar}\), ni un encadrement de Machin,
factoriel ou de Perron. Nous définissons la sous-relation en avant
\[
 \operatorname{Ext}^{\chi,\rightarrow}_{r+6,r}
 :=\{(x_r^\chi,\mathcal U_r^\chi(x_r^\chi))\}
 \subseteq\operatorname{Ext}_{r+6,r}.
\]
Ce n'est pas la sélection rétrospective
\(\operatorname{Ext}\cap\{P_{\rm ar}=\chi\}\) : le caractère est déjà une
coordonnée causale de \(x_{36}^\chi\). Les équations précédentes donnent
\(\lfloor N_{r+6}^\chi/729\rfloor=N_r^\chi\) et conservent feuillet, orientation,
retenue, bord, voie, signature, registre et généalogie.

La relation complète \(\operatorname{Ext}\) demeure multivoque.
L'unicité requise est celle de chaque section en avant sur son germe R36
produit,
\[
 \#\Bigl\{(x_r^\chi)_{r\ge36}:x_{36}^\chi\text{ fixé},\ 
 (x_r^\chi,x_{r+6}^\chi)
 \in\operatorname{Ext}^{\chi,\rightarrow}_{r+6,r}\Bigr\}=1,
\]
non la fausse égalité \(\#\operatorname{Ext}_{r+6,r}(x_r)=1\) à chaque porte
locale. Cette compatibilité se prolonge à tout niveau. La notation
\[
 w_6\longrightarrow w_{12}\longrightarrow w_{18}\longrightarrow w_{24}
 \longrightarrow w_{30}\longrightarrow R_{36}\longrightarrow G_9
\]
désigne ces applications et leurs changements de type, non une liste d'étiquettes.
L'holonomie \(\Gamma_9\) satisfait
\[
 g(k+9)=g(k),
 \qquad q^{(9)}(k+9)=q^{(9)}(k)+1,
\]
de sorte que la phase revient tandis qu'avancent le quotient, la voie et la
mémoire de l'état.

Une histoire canonique générique basée sur \(a_{(1,1)}\) et
\(t_j=M_{\rm ph}(\rho_9(j+1))\) atteste la non-vacuité de
\(\operatorname{Ext}\) ; elle n'individualise pas à elle seule les quatre sections
précédentes. Celles-ci sont raccordées par leurs coordonnées internes R36 et par la
règle en avant qui vient d'être définie.

\begin{theorem}[Construction inaugurale à partir des germes]
\label{thm:c27-construccion-inaugural-semillas}
Soit
\[
 I_9=\{0,\ldots,8\},\qquad D=\{N,E,S,O\},\qquad
 \mathcal S=I_9^2\times D,
\]
et soit \(\Omega_{\mathrm{seed}}=\mathcal S_+\times\mathcal S_\times\) le
domaine des deux curseurs qui réalisent les feuillets additif et multiplicatif
d'APP. Alors
\[
 |\Omega_{\mathrm{seed}}|=(9^2\cdot4)^2=104\,976,
\]
et la récurrence de \(54\) pas et la réduction ternaire définissent
\[
 \Omega_{\mathrm{seed}}
 \xrightarrow{\ T_{54}^{\rm seed}\ }
 \mathcal U_6^{\mathrm{cat}}
 \xrightarrow{\ \rho_3\ }
 \mathcal W_6
 \hookrightarrow\mathbb F_3^6,
\]
avec
\[
 |\mathcal U_6^{\mathrm{cat}}|=468,
 \qquad |\mathcal W_6|=243,
 \qquad |\mathbb F_3^6|=729.
\]
Les fibres satisfont exactement
\[
 432\cdot144+18\cdot432+18\cdot1944=104\,976.
\]
La relation d'extension du TPK prolonge les mots survivants à la
famille des sections compatibles de profondeur arbitraire. Les sous-relations
en avant ancrées dans R36 individualisent la génération coinductive
corrélée de \(\pi,\varphi,e\) et de \(\alpha\) ; la relation complète peut
continuer à posséder plusieurs prolongements locaux. Sur cette famille sont
construites conjointement l'ombre archimédienne et la réalisation
exceptionnelle du \cref{thm:c27-tesis-rectora-doble-proyeccion}.
\end{theorem}

\begin{proof}
Le produit cartésien précédent énumère exhaustivement les positions et
orientations des deux curseurs. L'application \(T_{54}^{\rm seed}\) est obtenue en
composant à chaque pas la sélection de feuillet, le transport cardinal, la
mise à jour du quotient et la mémoire de retenue. L'énumération de son
image contient \(468\) sextuplets distincts ; la somme des cardinalités de
ses fibres est l'identité exposée. L'application composante par composante
\(\rho_3\) contient \(243\) mots distincts. Enfin, la non-vacuité à
tout horizon et la compatibilité de \(\operatorname{Ext}\) avec les
troncatures produisent la limite inverse. La fonctionnalité de chaque
\(\operatorname{Ext}^{\chi,\rightarrow}\), à partir de l'état R36 qui
porte déjà \(\chi_{\rm HMT}\), démontre par récurrence l'unicité de cette section
globale sans exiger que chaque fibre locale de \(\operatorname{Ext}\) soit
un singleton. Chaque affirmation est une identité interne ou une conséquence de la
compatibilité projective ; aucune n'utilise en entrée un nombre réel, une
constante reconnue ou une table décimale.
\end{proof}

\begin{theorem}[Thèse directrice : un continu et deux projections]
\label{thm:c27-tesis-rectora-doble-proyeccion}
Soit \(\Omega_\Gamma^{\mathrm{cal}}\subset
X_\infty^{\mathrm{cal}}\) l'espace des sections survivantes des
tours enrichis du TPK. La retenue des trois tours
\(\gamma_-,\gamma_0,\gamma_+\) définit, sans évaluation archimédienne préalable,
l'homéomorphisme
\[
 \beta:\Omega_\Gamma^{\mathrm{cal}}\xrightarrow{\ \sim\ }
 (\mathbb F_3^6)^{\mathbb N}=: \mathcal W_\infty.
\]
Les formes équilibrées normalisées
\[
 a_k+c_k=\rho_k+3c_{k+1},\qquad
 \rho_k\in\{-1,0,+1\},
\]
construisent l'anneau ordonné \(\mathbb Z_{\rm HMT}\), son corps des
fractions \(\mathbb Q_{\rm HMT}\) et la complétion
\[
 \mathbb R_{\rm HMT}:=
 \operatorname{Cau}(\mathbb Q_{\rm HMT})/{\asymp}.
\]
La structure discrète conjointe est le produit fibré
\[
\mathcal C_{\mathrm{cont}}^{\mathrm{disc}}
=\bigl(\mathbb Z_{\rm HMT}\times
\Omega_\Gamma^{\mathrm{cal}}\bigr)
\times_{\mathcal W_\infty}X_\infty^{\mathrm{cal}}.
\]
Ses deux applications naturelles sont
\[
 P_{\mathrm{ar}}:\mathbb Z_{\rm HMT}\times
 \Omega_\Gamma^{\mathrm{cal}}\longrightarrow\mathbb R_{\rm HMT},
 \qquad
 Q_\infty:X_\infty^{\mathrm{cal}}\longrightarrow
 (\mathscr C_W^{\mathrm{cal}})^{\mathbb N}.
\]
Si \(\beta(\xi)=(b_q)_{q\geq0}\),
\(\nu(b)=\sum_{i=1}^{6}b_i3^{6-i}\) et
\[
 q_n(m,\xi)=m+\sum_{q=0}^{n-1}
 \frac{\nu(b_q)}{729^{q+1}},
\]
alors \((q_n)\) est de Cauchy et \(P_{\mathrm{ar}}(m,\xi)\) est sa classe
dans \(\mathbb R_{\rm HMT}\). Les cylindres rationnels
\[
 I_n(m;b_0,\ldots,b_{n-1})=
 \left[q_n,q_n+729^{-n}\right]
\]
sont compatibles, ont un diamètre tendant vers zéro et la partition récurrente
en \(729\) sous-cylindres construit, pour toute classe de Cauchy, un ruban
\((b_q)\) qui la représente. En conséquence,
\[
 P_{\mathrm{ar}}\ \text{est surjective},\qquad
 (\mathbb Z_{\rm HMT}\times\Omega_\Gamma^{\mathrm{cal}})
 /{\sim_{\rm ar}}\ \cong\ \mathbb R_{\rm HMT}.
\]
Ce n'est qu'après cette construction interne que l'unicité du corps ordonné
complet archimédien reconnaît \(\mathbb R_{\rm HMT}\cong\mathbb R\).

La seconde projection conserve l'incidence calibrée de chaque niveau ; dans les
coordonnées \(x=((w_j,f_j))_{j\geq1}\), elle agit par
\[
 Q_\infty(x)
 =\bigl(f_j,\Gamma_{f_jA_Wf_j^{-1}}(w_j)\bigr)_{j\geq1}.
\]
Ainsi, \(\mathbb R_{\rm HMT}\), reconnu ultérieurement comme \(\mathbb R\),
est l'ombre archimédienne obtenue en contractant la généalogie complète des
cylindres. La projection exceptionnelle conserve cette généalogie comme incidence
de bord et admet la réalisation successive
Paley--Witt--Golay--Leech--\(V^\natural\)--Monstre. Les deux branches proviennent
du même état et aucune n'est reconstruite rétrospectivement à partir de l'autre.
\end{theorem}

\begin{proof}
Les trois tours enrichis existent après chaque préfixe et leur retenue
les distingue par \(-1,0,+1\). La concaténation et la troncature sont
inverses, ce qui prouve la bijectivité et la continuité de \(\beta\). Pour une
classe de Cauchy de \(\mathbb Q_{\rm HMT}\), on choisit d'abord son cylindre
unitaire ordonné, puis, récursivement, l'unique sous-cylindre semi-ouvert
de la partition interne en \(729\) morceaux qui contient la classe ; sur un
bord commun, les deux écritures sont identifiées par \(\sim_{\rm ar}\).
La suite des extrémités gauches est précisément \((q_n)\), converge vers la
classe donnée et prouve la surjectivité sans présupposer un développement de
\([0,1]\) ni un point de \(\mathbb R\). L'identification au corps réel
est le théorème d'unicité appliqué ensuite à \(\mathbb R_{\rm HMT}\).

Les sous-espaces qui satisfont la fermeture d'un caractère précis sont des filtres
de survie au sein de la capacité globale. Leur fonction est d'individualiser
les sections de \(\pi,e,\varphi,\alpha\) ; la surjectivité du continu
provient de la complétion de \(\mathbb Q_{\rm HMT}\). Cette séparation des
domaines empêche
de transformer une condition de sélection de constantes en prémisse de la
couverture de \(\mathbb R\).

Dans la seconde branche, la compatibilité
\(r_{n,m}Q_n=Q_m p_{n,m}\) permet de passer à la limite et définit \(Q_\infty\).
La conjugaison par \(f_j\) transporte le calibre sans effacer l'incidence.
Les foncteurs ultérieurs de code, de réseau et d'algèbre de vertex réalisent
la chaîne exceptionnelle. L'égalité des rubans dans le produit fibré
construit la source conjointe et conserve des informations complémentaires : valeur
archimédienne dans la première ; incidence, bord et généalogie dans la seconde.
\end{proof}

\begin{corollary}[Développement infini comme objet mathématique]
\label{cor:c27-expansion-infinita-objeto}
Pour chaque caractère numérique \(\chi\) publié par l'état coinductif, la
lecture décimale est une séquence complète
\[
 \operatorname{Dec}_\chi(x_\infty)
 =(d_{\chi,1},d_{\chi,2},\ldots)
 \in\{0,\ldots,9\}^{\mathbb N}.
\]
Son préfixe de longueur \(N\) est la projection
\(\rho_N\operatorname{Dec}_\chi(x_\infty)\). Ainsi, mille, dix mille ou un
million de chiffres sont des coupes du même objet infini ; la génération ne
consiste pas à relever successivement une borne expérimentale.
\end{corollary}

\begin{theorem}[Génération coinductive corrélée de \(\pi\), \(\varphi\), \(e\) et \(\alpha\)]
\label{thm:c27-cuadrirrelacion-generada}
Le profil complet du système \(\mathfrak H_9\) détermine la famille
corrélée de sections en avant compatibles
\[
 \boldsymbol x_\infty=
 (x_\infty^\pi,x_\infty^\varphi,x_\infty^e,x_\infty^\alpha)
 \in
 \widehat\Omega_\pi\times_{R_{36}}
 \widehat\Omega_\varphi\times_{R_{36}}
 \widehat\Omega_e\times_{R_{36}}\widehat\Omega_\alpha,
 \qquad
 \widehat\Omega_\chi:=
 \varprojlim_{r\ge36}
 (\widehat{\mathcal X}^{\mathrm{enr},\chi}_r,
  \operatorname{Ext}^{\chi,\rightarrow}_{r+6,r}),
\]
et quatre caractères corrélés
\[
 \bigl(\Pi_{\mathrm{cl}},\Pi_{\mathrm{aut}},
       \Pi_{\mathrm{prop}},\Pi_{\mathrm{dod}}\bigr)
       (\boldsymbol x_\infty)
 =\bigl(\pi_{\mathrm{HMT}},\varphi_{\mathrm{HMT}},
        e_{\mathrm{HMT}},\alpha_{\mathrm{HMT}}\bigr).
\]
Les trois premiers caractères proviennent, respectivement, de l'unique
orbite de clôture, de l'unique orbite d'auto-échelle et de l'unique
orbite de propagation du catalogue enrichi. Le quatrième provient de la
coordonnée primitive signée de la famille terminale corrélée et de sa
fermeture dodécaphasique. Les quatre coordonnées sont donc des sorties de la
génération coinductive corrélée de profondeur arbitraire gouvernée par
la même dynamique APP--TRIT--TPK, avec retour de phase, avancée de mémoire et
conservation généalogique, bien que leurs applications de publication soient
distinctes.
\end{theorem}

\begin{proof}
Le recensement exhaustif
\(104,976\to468\to243\subset\mathbb F_3^6\) décompose l'image réalisée
en \(43\) orbites diédriques. Les invariants de fibre laissent une seule orbite
dans chacune des classes de clôture, de propagation et d'auto-échelle ; le calibre
interne y oriente, sans consultation décimale, les représentants
\[
 w_6^{\mathrm{cl}}=010211,
 \qquad w_6^{\mathrm{prop}}=201101,
 \qquad w_6^{\mathrm{aut}}=121200.
\]
Les sous-relations en avant conservent le caractère interne et les prédicats qui
définissent les trois classes ; Machin, la récurrence factorielle et Perron ne choisissent
aucun enfant de \(\operatorname{Ext}\). Dans le
canal de clôture, la composition à cinq régions avec retour unitaire produit
les rapports \(1/5\), \(1/239\) et la relation \(120/119\) ; l'identité de
Machin reconnaît exactement le demi-tour. Dans le canal de propagation, la
loi de semi-groupe normalisée impose
\((n+1)a_{n+1}=a_n\), donc \(a_n=1/n!\) et \(P_1=e\). Dans le canal
d'auto-échelle, l'incidence
\[
 F_{\mathrm{av}}=\begin{pmatrix}0&1\\1&1\end{pmatrix}
\]
possède un unique rayon positif et son équation propre impose
\(x^2-x-1=0\), \(x>0\), c'est-à-dire \(x=\varphi\).

Dans le profil TPK complet, la lecture terminale
\(x_{\mathrm{term}}\mapsto(B_{\mathrm{term}},U_{12}^{\mathrm{sgn}})\)
publie simultanément la carte des trois canaux et la coordonnée
primitive signée. La
seconde est identifiée réversiblement au sceau \(K\) et à ses différences
dodécaphasiques. La récurrence euclidienne avec retenue
\[
 s_m=P_m+E_m-\Phi_m-K_m+c_{m+1},\qquad
 A_m=s_m\bmod1000,\qquad
 c_m=(s_m-A_m)/1000
\]
produit \(\alpha_{\mathrm{HMT}}\) sans la recevoir en entrée. La
caractérisation analytique indépendante commence par le noyau fini
\(E_{21/22}^{(9)}\), qui détermine une racine simple \(\xi_9\) dans
\((0,10^{-2})\) sans consulter la sortie dodécaphasique. L'itération du même
lecteur gradué produit le système inverse
\(E_{21/22}^{(6)}\leftarrow E_{21/22}^{(9)}\leftarrow
E_{21/22}^{(12)}\leftarrow\cdots\). Ses troncatures commutent avec celles de la
voie dodécaphasique et les deux familles déterminent à chaque profondeur le même
cylindre survivant. Comme leurs diamètres tendent vers zéro, l'intersection est
un singleton et l'on obtient exactement
\(\alpha_A=\alpha_B=\alpha_{\mathrm{HMT}}\). L'égalité
\(\operatorname{Tr}_9(\xi_9^{-1})=
\operatorname{Tr}_9(\alpha_{12}^{-1})\) n'est que le certificat fini
d'ordre neuf de cette identité de sections complètes. Ainsi, la séparation
des lecteurs conserve l'origine commune des quatre coordonnées sans introduire
un prolongement en suspens ni une valeur cible.
\end{proof}

\begin{proposition}[Registre dodécaphasique de l'état complet antérieur à \(\alpha\)]
\label{prop:c27-registro-u-k-producido}
L'état terminal enrichi possède la double publication
\[
x_{\mathrm{term}}
\xrightarrow{(\pi_{\mathrm{term}},R_{\mathrm{sgn}})}
(B_{\mathrm{term}},U),
\]
où
\[
U=(2378,1406,2479,-452,998,-551,
-668,-204,-371,-322,-28,-997).
\]
En ordonnant les coordonnées selon les trois orbites de pas trois, la transformée
\[
\mathcal H_{12}
=P_{\mathrm{orb}}^{-1}(H_4\oplus H_4\oplus H_4)P_{\mathrm{orb}},
\qquad H_4^2=4I_4,
\]
possède l'inverse \(\mathcal H_{12}/4\) sur son image entière et détermine
de manière univoque
\[
K=(234,543,140,729,659,824,621,058,914,794,146,601).
\]
Par conséquent, au sein du profil TPK complet,
\[
x_{\mathrm{term}}\xrightarrow{R_{\mathrm{sgn}}}U
\xrightarrow{\mathcal H_{12}^{-1}}K
\]
est une composition exacte antérieure à la retenue qui publie \(\alpha\). La matrice
visible \(B_{\mathrm{term}}\) et le registre signé \(U\) sont deux
coordonnées du même état ; la seconde conserve les champs d'orientation,
d'opposition, de quotient et de registre que la première contracte. La composition
\[
 \Omega_{\mathrm{seed}}\longrightarrow
 \mathcal U_6^{\mathrm{cat}}\longrightarrow
 \widehat{\mathcal X}^{\mathrm{enr}}_{\mathrm{term}}
 \xrightarrow{R_{\mathrm{sgn}}}U
 \xrightarrow{\mathcal H_{12}^{-1}}K
\]
construit chaque coordonnée par APP--TRIT--TPK et conserve leurs registres. Ses
entrées sont les germes, les feuillets et les règles opératorielles ; \(U\), \(K\),
leurs différences, \(\alpha\) et les développements publiés appartiennent au
codomaine d'étapes successives.
\end{proposition}

\begin{theorem}[Clôture bidirectionnelle de la seconde voie de \(\alpha\)]
\label{thm:c27-alpha-segunda-via-bidireccional}
Le prolongement du noyau des lacunes jusqu'aux degrés \(7{:}9\) détermine
\[
\begin{aligned}
B_9(\alpha):={}&\log_{10}\!\left(\frac{10}{9}\right)
+\frac74\alpha^2-\frac{\alpha^4}{20}
+\frac{2\alpha^5}{21}+\frac{\alpha^6}{46}\\
&+\frac{\alpha^7}{120}-\frac{\alpha^8}{45}
-\frac{2\alpha^9}{495}.
\end{aligned}
\]
En relevant la rotation à l'inconnue positive \(T\), la résolvante équivaut à
\[
\boxed{\alpha T^2-B_9(\alpha)T+\alpha^3=0.}
\]
Avec
\(\xi=\operatorname{arcosh}(B_9(\alpha)/(2\alpha^2))\), les solutions sont
\[
T_\pm=\alpha e^{\pm\xi},
\qquad T_+T_-=\alpha^2,
\qquad
J_\alpha(T)=\frac{\alpha^2}{T},
\qquad J_\alpha(T_\pm)=T_\mp.
\]
La chambre \(6<T<7\) sélectionne \(T_+\) ; la clôture finie
\(\pi_9=T_+/2\) se prolonge au moyen de \(\Gamma_9\) jusqu'à
\(\pi_{\mathrm{HMT}}\). L'identité
\(\alpha=\sqrt{T_+T_-}\) fixe la moyenne géométrique des branches conjuguées.
Le jet \(7{:}9\), la racine simple, l'involution et le retour constituent la
seconde voie interne complète.
\end{theorem}

\begin{theorem}[Couplage exact du secteur \(\alpha\) avec le Monstre]
\label{thm:c27-alpha-rama-excepcional}
Dans la carte d'incidence du même registre dodécaphasique, les racines radiales
de \(A_2^{12}\) déterminent
\[
v_\alpha=4\rho_1+\rho_2+\cdots+\rho_{12},
\qquad v_\alpha^2=54,
\]
et le voisin d'indice trois
\[
N_{\alpha,0}
=\{x\in N(C_W):\langle x,v_\alpha\rangle\equiv0\pmod3\},
\qquad
N_\alpha=N_{\alpha,0}+\mathbb Z\frac{v_\alpha}{3}.
\]
Le réseau \(N_\alpha\) est pair, unimodulaire, de rang vingt-quatre et sans
racines, de sorte que \(N_\alpha\cong\Lambda_{24}\). La chaîne
\[
C_W\longrightarrow N(A_2^{12})
\xrightarrow{v_\alpha}\Lambda_{24}
\longrightarrow V^\natural
\xrightarrow{\operatorname{Aut}}\mathbb M
\]
expose la charnière exceptionnelle. Le scalaire \(\alpha_{\mathrm{HMT}}\) et le
vecteur \(v_\alpha\) appartiennent à des codomaines distincts et partagent l'état
dodécaphasique d'origine ; cette corrélation typée est le raccordement aux
symétries du Monstre.
\end{theorem}

\begin{proposition}[Statut des transports \(L_0,L_1\)]
\label{prop:c27-estatuto-l0-l1}
Les endomorphismes \(L_0,L_1\in\operatorname{End}(\mathbb F_3^6)\) et le
calendrier \((L_0,L_0,L_1,L_1)\) appartiennent à la loi de transport de l'état
TPK enrichi. Les six transitions publiées de rang plein fixent
chaque matrice de manière unique par \(L=X^{-1}Y\) dans \(\mathbb F_3\), et le
calendrier est le seul des \(16\) calendriers binaires qui conserve
simultanément les trois biographies structurelles. Sa fonction est de prolonger
les représentants déjà sélectionnés ; l'évaluation archimédienne et la
nomenclature conventionnelle sont appliquées ensuite.
\end{proposition}

Le domaine de cette proposition est \(\mathfrak H_9\), le système complet
APP--TRIT--TPK avec état enrichi. Le recensement APP constitue son étape
initiale ; TRIT apporte l'orientation et TPK construit le prolongement, la
mémoire et le retour. Cette séquence fixe exactement le domaine de chaque
flèche et l'emplacement de chaque sortie.

\begin{theorem}[Profondeur arbitraire et détermination des développements complets]
\label{thm:c27-profundidad-arbitraria}
Pour chaque
\(\chi\in\{\pi_{\mathrm{HMT}},e_{\mathrm{HMT}},
\varphi_{\mathrm{HMT}}\}\) et chaque \(N\ge1\), il existe un unique préfixe
décimal \(p_{\chi,N}\) de longueur \(N\), produit par une troncature
finie de l'état enrichi, tel que
\[
 \rho_{N+1,N}(p_{\chi,N+1})=p_{\chi,N}.
\]
La famille compatible est la famille des projections d'un unique élément
\[
 p_{\chi,\infty}
 \in\varprojlim_N\{0,\ldots,9\}^{N}
 \cong\{0,\ldots,9\}^{\mathbb N},
\]
et l'intersection de ses cylindres archimédiens est constituée d'un unique réel.
Par conséquent, la connexion engendre le développement complet ; chaque précision
finie est une projection de cette sortie infinie.
\end{theorem}

\begin{proof}
La preuve comporte deux étapes qui ne peuvent être inversées. Dans la première, entièrement
antérieure à la carte décimale, la section survivante de
\(\mathfrak H_9\) publie trois caractères exacts. Le caractère de clôture
est fixé par la pentafibre orientée, sa pente primitive \(1/5\), la
composition de ses quatre compagnes \(120/119\) et le compensateur entier
\(239\) ; le produit gaussien correspondant se termine sur la demi-droite
réelle négative et fixe de manière univoque le demi-tour orienté. Le caractère de
propagation est l'unique série formelle normalisée par
\(a_0=1\) et \((n+1)a_{n+1}=a_n\). Le caractère d'auto-échelle est l'unique rayon
expansif positif de
\[
 F_{\rm av}=\begin{pmatrix}0&1\\1&1\end{pmatrix},
 \qquad X^2-X-1=0.
\]
Ces trois conditions sont des sorties de l'état enrichi ; elles ne contiennent pas de
préfixes décimaux et ne consultent pas de valeurs conventionnelles.

Ce n'est qu'à la seconde étape qu'agit la publication archimédienne. Pour chaque caractère
déjà généré et chaque \(N\), soit \(\operatorname{Dec}_N\) l'unique cylindre
décimal canonique de diamètre \(10^{-N}\) qui contient sa réalisation exacte.
Ce cylindre est une coordonnée publiée du caractère : ce n'est pas une fibre locale
de \(\operatorname{Ext}\) et il n'est pas utilisé pour la résoudre.
L'unicité de la réalisation et la convention non terminale impliquent
\[
 \rho_{N+1,N}\operatorname{Dec}_{N+1}(\chi)
 =\operatorname{Dec}_N(\chi).
\]
Les identités de Machin, les sommes partielles de la série exponentielle et les
quotients consécutifs de Fibonacci--Perron fournissent des intervalles rationnels
certifiés pour calculer \(\operatorname{Dec}_N\) ; ils n'interviennent pas dans
\(\operatorname{Ext}\), ne choisissent ni orbite ni caractère et ne sélectionnent pas
d'enfants de ses fibres locales. Ils calculent bien le bloc de coordonnées de publication,
qui demeure typé comme sortie ultérieure. Les
cylindres publiés sont emboîtés et leurs diamètres tendent vers zéro, de sorte
que leur limite inverse détermine un unique point. Comme \(N\) est arbitraire et
que la relation structurelle d'extension est dépourvue de profondeur terminale, la
conclusion vaut simultanément pour toute longueur finie et pour le développement
infini.
\end{proof}

Le porteur exécutable impose matériellement cette séparation au moyen de trois
fonctions disjointes : il engendre d'abord les caractères exacts, puis les évalue
et ne publie qu'à la fin les blocs archimédiens. Son propre audit syntaxique
échoue si un évaluateur entre dans le générateur. L'exécution de \(396\) niveaux,
\(2376\) trits, \(43\) tours, \(387\) couples \((K,K+9)\) par canal et
\(1000\) chiffres vérifie une vaste instance matérielle du théorème ; celle de
\(10,000\) chiffres en vérifie une autre. Aucune de ces deux quantités ne limite la
profondeur. Le quantificateur \(\forall N\) et la compatibilité
\(\rho_{N+1,N}p_{N+1}=p_N\), avec l'existence de l'élément
\(p_\infty\) de la limite inverse, démontrent la sortie décimale infinie ; les
calculs finis sont des certificats reproductibles de ses projections.

La coordonnée \(\alpha_{\mathrm{HMT}}\) appartient à la même famille limite
corrélée. Sa voie entière se prolonge au moyen de
\(\mathsf T_{\alpha,\infty}\). Le lecteur
analytique interne construit exactement \(E_{21/22}^{(9)}\) et sa racine simple
\(\xi_9\) dans le quotient des jets d'ordre neuf. La stabilité vérifiée à
\(10,000\) chiffres certifie la propagation distante de la retenue de la voie
entière, mais n'est qu'une projection finie. Les carrés de troncature de
la famille dodécaphasique et de la famille de jets commutent exactement ;
l'unicité de leurs sections compatibles identifie
\(\alpha_A=\alpha_B=\alpha_{\mathrm{HMT}}\) dans
\(\mathbb R_{\rm HMT}\) avant toute comparaison métrologique. Les exécutions à
neuf, mille ou dix mille chiffres sont des témoins finis de cette identité et non sa
portée.

\begin{proposition}[Prolongement complet de la voie entière de \(\alpha\)]
\label{prop:c27-alpha-prolongacion-completa}
Soit \(B=1000\), soient \((P_k),(E_k),(\Phi_k)\) les trois suites de blocs
publiées par l'état coinductif et soit \((K_k)\) la suite
dodécaphasique prolongée. Les quatre séries
\[
 p=\sum_{k\geq1}P_kB^{-k},\qquad
 e_0=\sum_{k\geq1}E_kB^{-k},\qquad
 \phi_0=\sum_{k\geq1}\Phi_kB^{-k},\qquad
 \kappa=\sum_{k\geq1}K_kB^{-k}
\]
convergent absolument et déterminent
\[
 \alpha_A=p+e_0-\phi_0-\kappa.
\]
La récurrence de retenue est l'écriture en base \(B\) de cette égalité.
Pour chaque \(M\), une troncature suffisamment profonde détermine les
premiers \(M\) chiffres du développement canonique de \(\alpha_A\) ; l'erreur de
queue est bornée par une constante multipliée par \(B^{-N}\). Par
conséquent, \(\mathsf T_{\alpha,\infty}\) publie une suite décimale
complète, et ses exécutions à \(3000\) ou \(10\,000\) chiffres sont des projections
finies de cette sortie.
\end{proposition}

\begin{proof}
Chaque bloc appartient à un ensemble entier borné, de sorte que les quatre
séries sont dominées par une série géométrique. La somme partielle jusqu'à \(N\)
diffère de la somme complète de \(O(B^{-N})\). La division euclidienne de
droite à gauche transporte exactement cette queue au moyen de la retenue ;
par unicité du développement canonique, tout préfixe qui reste séparé d'un
bord de cylindre se stabilise lorsque \(N\) augmente, et sur un bord
on applique la convention canonique non terminale. Cela définit simultanément
tous les chiffres, non une suite d'objectifs numériques indépendants.
\end{proof}

\begin{corollary}[Séparation entre génération et reconnaissance]
\label{cor:c27-generacion-reconocimiento}
Les identités de Machin, la série exponentielle, le rayon de Perron et la carte
décimale sont des réalisations exactes de caractères préalablement sélectionnés
par \(\mathfrak H_9\). Les noms conventionnels \(\pi,e,\varphi,\alpha\)
et les tables métrologiques reconnaissent les sorties ; ils ne déterminent ni orbites,
ni caractères, ni transports, ni retenues. Les préfixes qu'ils calculent sont uniquement
des projections coordonnées ultérieures des caractères déjà générés.
\end{corollary}

\paragraph{Évaluation des caractères construits.}
À l'étape d'évaluation archimédienne, le caractère de propagation
détermine la condition initiale \(a_0=1\) et la récurrence
\((n+1)a_{n+1}=a_n\). Le calcul des termes applique cette récurrence.
Pour l'auto-échelle, le vecteur ligne évolue au moyen de l'incidence
préalablement construite,
\[
 (x_{n+1},y_{n+1})=(x_n,y_n)F_{\rm av},
 \qquad
 F_{\rm av}=\begin{pmatrix}0&1\\1&1\end{pmatrix},
 \qquad
 \det(tI-F_{\rm av})=t^2-t-1.
\]
Les quotients de composantes fournissent les bornes rationnelles du
lecteur d'auto-échelle. L'implémentation vérifie la concordance entre
incidence, polynôme caractéristique, orientation et normalisation avant
d'effectuer cette évaluation. Une altération ou une absence de ces données
interrompt le calcul au lieu de conserver tacitement la valeur
précédente. Ce contrôle relie les caractères déjà produits à leurs
évaluateurs ; sa place dans la construction est postérieure au générateur.
